\chapter{Thiết kế hệ thống}

\section{Kiến trúc tổng thể}

Hệ thống ShopNexus được thiết kế theo kiến trúc microservices, trong đó mỗi
thành phần đảm nhận một chức năng riêng biệt và giao tiếp với nhau thông qua
các giao thức TCP và HTTP. Việc tách biệt các dịch vụ cho phép phát triển,
triển khai và mở rộng từng thành phần một cách độc lập.

Toàn bộ hệ thống bao gồm 5 thành phần chính. Thiết bị nhúng ESP32-CAM đóng
vai trò đầu vào, thu âm thanh qua microphone INMP441 và hình ảnh qua camera
OV5640, sau đó truyền dữ liệu real-time đến server thông qua TCP. Server
trung gian chạy trên Node.js hoạt động như cầu nối giữa thiết bị nhúng và
trình duyệt, thực hiện relay TCP$\leftrightarrow$WebSocket đồng thời đóng vai
trò proxy chuyển tiếp request đến các dịch vụ AI phía sau. Ba dịch vụ AI
được triển khai riêng biệt: dịch vụ khử nhiễu sử dụng mô hình DTLN trên Flask,
dịch vụ nhận dạng giọng nói và sinh mô tả sản phẩm trên FastAPI, và dịch vụ
phân loại sản phẩm sử dụng XLM-RoBERTa cũng trên FastAPI.
\Cref{fig:system-architecture} minh họa kiến trúc tổng thể của hệ thống.

\begin{figure}[htbp]
\centering
\fbox{\parbox{0.85\linewidth}{\centering\vspace{2.5cm}
\textit{Thêm hình: Sơ đồ kiến trúc tổng thể hệ thống --- thể hiện 5 thành phần
(ESP32-CAM, Node.js Server, Noise Reduction, Voice Rephraze, Product Classification)
cùng các kết nối TCP/HTTP/WebSocket giữa chúng và trình duyệt.}
\vspace{2.5cm}}}
\caption{Kiến trúc tổng thể hệ thống ShopNexus}
\label{fig:system-architecture}
\end{figure}

\begin{table}[htbp]
\centering
\caption{Tổng quan các thành phần hệ thống}
\label{tab:system-overview}
\begin{tabular}{>{\raggedright\arraybackslash}p{0.22\linewidth} >{\centering\arraybackslash}p{0.18\linewidth} >{\centering\arraybackslash}p{0.18\linewidth} >{\raggedright\arraybackslash}p{0.3\linewidth}}
\toprule
\rowcolor{headerblue} \textcolor{white}{\textbf{Thành phần}} & \textcolor{white}{\textbf{Công nghệ}} & \textcolor{white}{\textbf{Port}} & \textcolor{white}{\textbf{Vai trò}} \\
\midrule
\rowcolor{rowgray} ESP32-CAM & Arduino/C++ & TCP 3001, 3002 & Thu âm thanh + hình ảnh \\
Server trung gian & Node.js + Express & 3000 & Relay TCP$\leftrightarrow$WebSocket, proxy API \\
\rowcolor{rowgray} Noise Reduction & Flask + TensorFlow & 5000 & Khử nhiễu DTLN \\
Voice Rephraze & FastAPI + OpenAI & 5001 & STT + sinh mô tả \\
\rowcolor{rowgray} Product Classification & FastAPI + PyTorch & 5002 & Phân loại 32 danh mục \\
\bottomrule
\end{tabular}
\end{table}

\section{Luồng xử lý chính}

Dữ liệu trong hệ thống đi qua một pipeline tuần tự gồm 5 bước. Đầu tiên,
ESP32 thu âm thanh từ microphone INMP441 ở định dạng 16-bit PCM, tần số lấy
mẫu 16 kHz, mono và truyền qua TCP port 3002 đến server trung gian. Tại đây,
server Node.js nhận stream PCM, đóng gói thành file WAV rồi gửi đến dịch vụ
khử nhiễu (port 5000) để loại bỏ tạp âm. Audio đã được làm sạch tiếp tục
được chuyển đến dịch vụ Voice Rephraze (port 5001), nơi giọng nói được nhận
dạng thành văn bản và tự động sinh mô tả sản phẩm theo phong cách viết do
người dùng lựa chọn. Văn bản mô tả sau đó được gửi đến dịch vụ phân loại
(port 5002) để xác định danh mục sản phẩm phù hợp trong 32 danh mục thương
mại điện tử. Cuối cùng, toàn bộ kết quả --- bao gồm spectrogram trước/sau
khử nhiễu, text nhận dạng, mô tả đã sinh và top-K danh mục dự đoán ---
được hiển thị trên giao diện web.
\Cref{fig:data-pipeline} minh họa luồng xử lý này.

\begin{figure}[htbp]
\centering
\fbox{\parbox{0.85\linewidth}{\centering\vspace{2cm}
\textit{Thêm hình: Sơ đồ luồng xử lý dữ liệu (flowchart hoặc sequence diagram)
--- từ ESP32 thu âm $\rightarrow$ Node.js đóng gói WAV $\rightarrow$ Denoise (port 5000)
$\rightarrow$ STT + sinh mô tả (port 5001) $\rightarrow$ Phân loại (port 5002)
$\rightarrow$ Hiển thị kết quả trên web UI.}
\vspace{2cm}}}
\caption{Luồng xử lý dữ liệu trong hệ thống}
\label{fig:data-pipeline}
\end{figure}

\section{Thiết kế thiết bị nhúng ESP32}

\subsection{Kiến trúc Dual-Core}

Một trong những quyết định thiết kế quan trọng nhất trên ESP32 là việc phân
tách tác vụ thu âm thanh và thu hình ảnh lên hai nhân xử lý riêng biệt.
Core 0 được gán riêng cho Audio Capture Task --- đọc I2S liên tục từ INMP441,
chuyển đổi từ 32-bit sang 16-bit PCM và stream qua TCP port 3002 với stack
16 KB. Core 1 đảm nhận Camera và Main Loop --- chụp JPEG từ OV5640, gửi qua
TCP port 3001, đồng thời xử lý các lệnh điều khiển camera từ server.

\begin{table}[htbp]
\centering
\caption{Phân công xử lý dual-core trên ESP32}
\label{tab:dual-core}
\begin{tabular}{>{\centering\arraybackslash}p{0.12\linewidth} >{\raggedright\arraybackslash}p{0.25\linewidth} >{\raggedright\arraybackslash}p{0.5\linewidth}}
\toprule
\rowcolor{headerblue} \textcolor{white}{\textbf{Core}} & \textcolor{white}{\textbf{Nhiệm vụ}} & \textcolor{white}{\textbf{Chi tiết}} \\
\midrule
\rowcolor{rowgray} Core 0 & Audio Capture Task & Đọc I2S liên tục từ INMP441, chuyển đổi 32-bit$\rightarrow$16-bit PCM, stream qua TCP port 3002. Stack: 16 KB. \\
Core 1 & Camera + Main Loop & Chụp JPEG từ OV5640, gửi qua TCP port 3001, xử lý lệnh điều khiển camera từ server. \\
\bottomrule
\end{tabular}
\end{table}

Việc tách riêng hai core đảm bảo luồng audio chạy liên tục không bị gián đoạn
bởi quá trình chụp ảnh từ camera, tránh hiện tượng buffer underrun --- một vấn
đề phổ biến khi xử lý đồng thời nhiều nguồn dữ liệu trên vi điều khiển.

\begin{figure}[htbp]
\centering
\fbox{\parbox{0.85\linewidth}{\centering\vspace{2cm}
\textit{Thêm hình: Sơ đồ kiến trúc dual-core ESP32 --- Core 0 chạy Audio Task
(I2S INMP441 $\rightarrow$ TCP 3002), Core 1 chạy Camera Task
(OV5640 $\rightarrow$ TCP 3001) + xử lý lệnh điều khiển.
Thể hiện rõ sự phân tách tác vụ giữa hai nhân.}
\vspace{2cm}}}
\caption{Kiến trúc dual-core trên ESP32}
\label{fig:dual-core}
\end{figure}

\subsection{Cấu hình chân kết nối}

Camera OV5640 sử dụng bus dữ liệu 8-bit song song với các chân D0--D7 được ánh
xạ tới GPIO 5, 18, 19, 21, 36, 39, 34, 35. Các tín hiệu điều khiển bao gồm
XCLK (GPIO 0), PCLK (GPIO 22), VSYNC (GPIO 25), HREF (GPIO 23), và giao tiếp
I2C qua SDA (GPIO 26), SCL (GPIO 27). Chân PWDN (GPIO 32) dùng để điều khiển
chế độ tiết kiệm năng lượng.

Microphone INMP441 kết nối qua giao thức I2S với 3 chân: WS -- Word Select
(GPIO 2), SCK -- Serial Clock (GPIO 14), và SD -- Serial Data (GPIO 15),
sử dụng I2S port 1 để tránh xung đột với các ngoại vi khác trên port 0.

\begin{figure}[htbp]
\centering
\fbox{\parbox{0.85\linewidth}{\centering\vspace{2.5cm}
\textit{Thêm hình: Sơ đồ đấu nối phần cứng (wiring diagram) --- ESP32-CAM
kết nối với camera OV5640 (bus 8-bit, I2C) và microphone INMP441 (I2S: WS, SCK, SD).
Ghi rõ số GPIO cho từng chân.}
\vspace{2.5cm}}}
\caption{Sơ đồ kết nối phần cứng ESP32-CAM}
\label{fig:wiring-diagram}
\end{figure}

\designDenoise
\designVoice
\designClassify

\section{Thiết kế server trung gian}

Server trung gian là thành phần kết nối then chốt trong kiến trúc hệ thống,
chạy trên Node.js/Express tại port 3000. Thành phần này giải quyết sự khác
biệt về giao thức giữa ESP32 (TCP) và trình duyệt (WebSocket/HTTP), đồng
thời đóng vai trò điều phối request đến các dịch vụ AI.

Về phía thiết bị nhúng, server duy trì hai TCP server: TCP Camera Server
(port 3001) nhận JPEG frame từ ESP32, parse header 4 byte chứa độ dài frame
rồi broadcast qua WebSocket endpoint \texttt{/camera} đến trình duyệt;
TCP Audio Server (port 3002) nhận stream 16-bit PCM, duy trì byte-alignment
và broadcast qua WebSocket endpoint \texttt{/audio}.

Về phía trình duyệt, server cung cấp HTTP API Proxy chuyển tiếp request đến
các dịch vụ AI: \texttt{POST /api/process-voice} gọi tuần tự Denoise Service
rồi STT Service, \texttt{POST /api/gen-description} gọi Voice Rephraze Service,
và \texttt{POST /api/classify} gọi Classification Service.
\Cref{fig:middleware-architecture} minh họa kiến trúc server trung gian.

\begin{figure}[htbp]
\centering
\fbox{\parbox{0.85\linewidth}{\centering\vspace{2.5cm}
\textit{Thêm hình: Sơ đồ kiến trúc server trung gian --- thể hiện ESP32
kết nối TCP (port 3001, 3002) đến Node.js, Node.js relay qua WebSocket
(/camera, /audio) đến Browser, đồng thời proxy HTTP API đến 3 dịch vụ AI.}
\vspace{2.5cm}}}
\caption{Kiến trúc server trung gian}
\label{fig:middleware-architecture}
\end{figure}

Giao diện web tích hợp hiển thị video streaming real-time, biểu đồ amplitude
âm thanh 30 bars, spectrogram trước và sau khử nhiễu thông qua thư viện
WaveSurfer.js, cùng kết quả nhận dạng và phân loại sản phẩm.

\begin{figure}[htbp]
\centering
\fbox{\parbox{0.85\linewidth}{\centering\vspace{2.5cm}
\textit{Thêm hình: Ảnh chụp màn hình giao diện web --- hiển thị các thành phần:
video streaming từ camera, biểu đồ amplitude âm thanh, spectrogram trước/sau
khử nhiễu, kết quả nhận dạng giọng nói, mô tả sản phẩm đã sinh,
và danh sách top-K danh mục dự đoán.}
\vspace{2.5cm}}}
\caption{Giao diện web hệ thống ShopNexus}
\label{fig:web-ui}
\end{figure}
